\chapter{Conclusions and Future Work}
\label{ch:conclusions}

\todo{AUDIT (MAJOR FIB defect): No \textbf{technical-competences map} is provided. The FIB tribunal rubric requires that the GEI/Computacio technical competences declared in the TFG inscription (e.g.\ CCO1.1, CCO2.1, CES1.x, CT codes) be explicitly mapped to the chapters/sections where each is demonstrated. Add a dedicated \texttt{\textbackslash section\{Technical Competences Achieved\}} with a table (competence code -> short name -> evidence section/figure). Without this, the tribunal will dock the rubric explicitly. Cross-check with the competences declared in your TFG inscription form.}

\todo{AUDIT (chapter length): At ~169 lines the chapter risks reading as a one-page recap. Target 4--6 pages by adding: (i) the competences map above, (ii) a Limitations subsection (currently limitations are scattered through paragraphs 5--6 of \S\ref{sec:conclusions}, not enumerated), (iii) a brief Methodological Reflection (~half a page: what was learned about iterative experimentation, dataset-driven debugging, calibration of confidence gates, etc.).}

\todo{AUDIT (missing section): No explicit \textbf{Limitations} subsection. The current narrative mentions scope limits in passing but the tribunal expects an itemised list (e.g.\ monophonic only, treble only, single-page only, no real-data ground truth, CTC-gate threshold checkpoint-bound, EasyOCR chord vocabulary, etc.). Add \texttt{\textbackslash subsection\{Limitations\}} with 5--8 bullets.}

\todo{AUDIT (missing section): No \textbf{Methodological Reflection}. Add ~half a page describing the iterative design loop (synthetic-first, calibrate-then-evaluate, decouple-melody-from-chord), what would be done differently in hindsight, and how the GEP/PM artefacts (Gantt, risk register) tracked against reality.}

\section{Conclusions}
\label{sec:conclusions}

This thesis has presented the design, implementation, and quantitative
evaluation of an end-to-end Optical Music Recognition system tailored
to the monophonic jazz lead sheet, with \textit{The Real Book} as the
primary target.
The system is organised as a \term{two-stream} architecture
(\Cref{sec:design-twostream}): a melody branch that consumes a
horizontal staff strip through a ResNet18 convolutional backbone,
a two-layer bidirectional LSTM with hidden size~256, and a CTC
output head, producing a sequence of LMX tokens that jointly
encode pitch, duration, accidentals, and barline structure
(\Cref{sec:lmx-format}); and a chord branch that runs an
architecturally identical CRNN--CTC model on the chord-symbol
strip extracted above each staff.
The two models are trained independently on the synthetic
LilyJAZZ-rendered PrIMuS corpus augmented with a calibrated
scan-simulation pipeline, and are wired together behind a single
FastAPI endpoint that accepts a PDF or image of a lead-sheet page
and returns a structured JSON transcription suitable for downstream
musical applications (\Cref{ch:design}, \Cref{ch:impl}).

On the headline figures, the music CRNN reaches an aggregate
Symbol Error Rate of \SI{1.19}{\percent} on the clean test split
and \SI{1.28}{\percent} on the scan-augmented split, with a
melodic SER---computed after stripping the structural \code{measure}
and tie tokens---of \SI{0.11}{\percent} and \SI{0.18}{\percent}
respectively, and a perfect-transcription rate of
\SI{73}{\percent} on the clean split and \SI{71}{\percent} on
the scanned split (\Cref{tab:results-headline}).
These numbers, obtained from the epoch-37 checkpoint on \num{4604}
held-out samples per split, demonstrate that a moderately sized
CRNN--CTC stack (\Cref{sec:soa-crnn}) is sufficient for practical
monophonic jazz transcription, and that the scan-simulation
augmentation strategy designed in \Cref{sec:design-data} narrows
the synthetic-to-real domain gap without requiring real annotated
training data: the gap between clean and scanned SER is
approximately \num{0.09} percentage points in aggregate and
\num{0.07} percentage points on the melodic metric.

Reviewing the specific objectives listed in \Cref{sec:objectives}:
Objective~1 (state-of-the-art survey and architectural
justification) is met by \Cref{ch:soa} and the comparative
analysis in \Cref{ch:design}.
Objective~2 (synthetic dataset construction via LilyPond and
LilyJAZZ, paired with LMX annotations) is met by the data
pipeline described in \Cref{sec:design-data}.
Objective~3 (scan-simulation augmentation) is met and validated
quantitatively by the small clean-vs-scanned gap reported above.
Objective~4 (CRNN--CTC design, training, and SER evaluation) is
met by the model and training stack of \Cref{ch:impl} and the
results of \Cref{sec:res-ser}.
Objective~5 (quantitative error analysis) is met by the per-token
breakdown of \Cref{sec:res-errors}, which decomposes the residual
SER into substitution, insertion, and deletion classes.
Objective~6 (complete inference pipeline with chord OCR returning
structured JSON) is met by the FastAPI service.\todo{AUDIT: ``met by the FastAPI service'' is too brief and lacks a \texttt{\textbackslash Cref} to the implementation section/endpoint description (e.g.\ \Cref{ch:impl} or \texttt{sec:impl-api}). Add the cross-reference so the claim is auditable.}
Objective~7 (domain-gap quantification and fine-tuning on real
Real Book pages) is \textit{partially} met: the synthetic
domain gap is bounded by the scanned-split evaluation, but the
fine-tuning stage on a manually annotated set of Real Book pages
remains an open extension (\Cref{sec:res-domaingap}).\todo{AUDIT (cross-chapter inconsistency): \S\ref{sec:res-domaingap} in Ch.~6 is currently empty (three unresolved \texttt{\textbackslash todo} placeholders, no figures, no numbers). Claiming Objective~7 is ``partially met'' while citing that section is unsupported. Either (a) execute the experiments and populate \S\ref{sec:res-domaingap}, or (b) downgrade this claim to ``not met / deferred to future work'' and drop the \texttt{\textbackslash Cref}. Tribunal will follow the link and find nothing.}

\todo{AUDIT (critical assessment is shallow): The four ``lessons learned'' below are positive findings, not a \textbf{critical} assessment. Add a paragraph that openly discusses what did NOT work or remains unsatisfying: e.g.\ residual scan-domain gap not measured on genuine pages, no comparison to commercial OMR (Audiveris, SmartScore), chord OCR via EasyOCR is a stopgap rather than a learned solution, CTC-confidence reject gate is checkpoint-bound and must be recalibrated. The tribunal rewards self-criticism.}
Four lessons learned are worth recording for follow-up work.
\textit{First}, the LMX token format is decisively more practical
than full MusicXML for staff-level sequence-to-sequence modelling:
it is compact, CTC-friendly, and avoids the deep nesting that
makes MusicXML cumbersome as a model target.
\textit{Second}, CTC remains a strong baseline for OMR despite
the popularity of transformer decoders: it requires no
autoregressive decoding loop, handles variable-length staves
without explicit segmentation, and trains stably on the data
volumes available within a Bachelor's thesis budget.
\textit{Third}, scan-simulation augmentation in
\code{albumentations} closes most of the synthetic-to-real
distribution shift, but residual errors persist on difficult
fonts and low-contrast pages---scanned-split SER never quite
collapses to the clean-split SER.
\textit{Fourth}, decoupling melody recognition and chord
recognition into two independent CRNNs prevents vocabulary
pollution between the two visual languages and allows each
branch to be retrained or fine-tuned without disturbing the
other; this modularity proved more valuable in practice than the
theoretical elegance of a unified output head would have been.

The scope of the system remains deliberately narrow
(\Cref{sec:scope}): monophonic melodies in treble clef only, on
the eight in-scope time signatures, with single-page PDFs and
no tuplets, voltas, or navigation marks.
Polyphonic writing, bass clef, handwritten manuscripts, and
multi-page layout recovery are explicitly out of scope.
The most important open question---visible in
\Cref{sec:res-domaingap}---is the model's accuracy on genuine
scanned Real Book pages, which the present evaluation bounds via
the scan-augmented synthetic split but does not measure end to
end against a manually annotated real-page ground truth.
Closing that gap is the natural next step and is the leading
item in the future-work plan that follows.

\section{Future Lines of Research}
\label{sec:future-work}

\begin{itemize}
  \item \textbf{Polyphonic OMR.}\todo{AUDIT (future-work concreteness): Future-work items are well-written but only one (real-data fine-tuning) gives a numeric target. Add a concrete scoping line to each remaining item: e.g.\ for Polyphonic, ``annotate N polyphonic staves from corpus X and target SER $\le$ Y\,\%''; for Full-page segmentation, ``train YOLOv8-n on the IMSLP-layout subset of MUSCIMA++ with $\ge$ 1k pages''; for Transformer decoder, ``reproduce SMT on the present LMX vocab with batch=8 on a 12\,GB GPU''. ``Explore X'' is vague; ``fine-tune model M on dataset D of size N'' is what the rubric rewards.}
        The current label grammar and CTC loss assume a single
        voice per staff.
        Extending the system to polyphonic input would require a
        richer LMX vocabulary (chord-stack tokens, voice
        indicators) and either a multi-target CTC formulation or
        a sequence-to-sequence decoder.
        The main challenge is label ambiguity in vertically
        stacked noteheads, where the reading order is not always
        recoverable from the image alone.

  \item \textbf{Real-data fine-tuning.}
        Collecting and manually annotating \numrange{50}{100}
        staves from genuine Real Book scans and fine-tuning both
        CRNN branches on this small set would directly close the
        residual synthetic-to-real gap identified in
        \Cref{sec:res-domaingap}.
        Standard transfer-learning practice suggests freezing the
        convolutional backbone and updating only the BiLSTM and
        CTC head, which limits the data requirement and reduces
        the risk of overfitting to the small fine-tuning set.

  \item \textbf{Full-page segmentation.}
        The current preprocessing stage relies on a
        morphology-and-projection staff detector that performs
        well on cleanly scanned single-page PDFs but is fragile
        on slanted staves, partial systems, and multi-page
        documents.
        Replacing it with a deep-learning detector (YOLO-style
        box regression or a DETR-style transformer detector
        trained on annotated page-layout data) is the most
        pressing engineering bottleneck for practical
        deployment and would unlock multi-page Real Book
        processing.

  \item \textbf{Transformer-based decoder.}
        The greedy CTC decoder currently used could be compared
        head-to-head against an autoregressive transformer
        decoder conditioned on the same convolutional features,
        in the style of the Sheet Music
        Transformer~\cite{riosVilaSheetMusicTransformer2024}.
        A transformer decoder would be expected to handle
        long-range dependencies between measures (key-signature
        propagation, beam-internal rhythm consistency) more
        gracefully than CTC, at the cost of slower inference
        and a larger data requirement; quantifying that trade-off
        on the present dataset is a natural follow-up
        experiment.

  \item \textbf{Richer output format (MusicXML / LMX v2).}
        The current LMX output is purpose-built for CTC training
        and omits notation features---ties across barlines, slurs,
        articulations, dynamics, double accidentals---that are
        present in many Real Book pages but not in the in-scope
        vocabulary.
        Exporting recognised staves to MusicXML would enable
        round-trip compatibility with notation software
        (MuseScore, Finale, Sibelius) and would require a
        post-processing pass that infers the missing
        notation-graphical structure from the LMX token stream
        and the original image.
\end{itemize}
